\cleardoublepage
\thispagestyle{empty}

\begin{center}
  \bfseries
  {\selectlanguage{english}Abstract}
\end{center}
\normalfont

\vspace{0.5cm}

This thesis studies automatic \kvp{} extraction from business documents with generative and layout-aware discriminative methods. The discriminative system combines a LayoutLMv3-base encoder, a token classifier, and a span-level linker. It uses extracted text and bounding boxes with constant blank visual input, and therefore receives no document-specific image information.

\textbf{Background:} Business document processing remains a significant challenge in document understanding, requiring accurate extraction of structured information from semi-structured layouts. Traditional approaches rely on template matching or rule-based systems, limiting generalization to diverse document types. Recent generative baselines using large language models achieve strong text extraction, but spatial localization can be less accurate and generated outputs are not constrained to observed input tokens.

\textbf{Method:} The two-head architecture assigns each token a role (key, value, or other) and links predicted key spans to value spans. We investigate token-level linking (V1), span-level linking (V2), validation-based diagnosis and corrections, and a final corrected configuration (V4). Because source PDFs were unavailable, broken, or lacked a usable text layer, the prepared development set contains 5{,}389 of the 9{,}656 unique KVP10k training pages, which are then split into training and validation subsets.

\textbf{Evaluation:} We follow IBM's released KVP10k benchmark: type-aware one-to-one pair matching with NED $<0.2$, IoU $\geq0.3$ in the executable code, and macro-averaged per-page precision and recall. Entity extraction and key-value linking are evaluated separately.

\textbf{Results:} V4 checkpoint selection uses validation Regular text+location macro F1 and selects epoch 10. On the test set, V4 obtains Regular F1 values of 0.392 for text, 0.446 for location, and 0.345 for text+location. Corrected class-aware entity micro and macro F1 are 0.792 and 0.772. The reconstructed Mistral baseline obtains 0.521 text+location F1. An unchanged recovery method raises V4 All text+location F1 from 0.229 to 0.261 without changing Regular F1. V4 reaches 0.305 and 0.410 Regular text+location F1 in the two annotation-geometry clusters. The numerically higher entity scores, together with qualitative errors, suggest that relation linking remains a major limitation of the current pipeline.

\textbf{Keywords:} Document Understanding, Key-Value Pair Extraction, LayoutLMv3, Span-Level Linking, Layout-Aware Learning, Information Extraction

\clearpage
